\section{Role integrity and the \ribench{} scenarios}\label{sec:benchmark}

\subsection{Obligations, interventions and endpoints}\label{sec:design}

A \emph{role obligation} names a use the senior role makes of
subordinate material and the evidential constraint that use must
respect. The three we study, to weigh, to verify and to relay, are
defined with their scenarios. Measuring
whether an obligation was met cannot
rely on agreement with a reference outcome, because a chair that counts
votes and one that weighs arguments can reach the same decision from
the same input. We therefore construct the subordinate material so that
the obligation and the easier behavior pull in different directions,
hold everything else fixed, and read the senior role's own decision,
ruling or report against that input. The cases come from real
institutions whose rules assign the obligation to the senior role and
whose records are long enough to make it hard to meet. Following these
principles, we score each episode from the senior agent's own output,
combining rule-based checks written alongside each case with a
rubric-guided LLM judge.

Each scenario reports one endpoint, a rate over a stated set of scored
units, together with its change from the Vanilla to the Obligation
prompt. The units and denominators differ by
scenario and are stated with each readout below. One minus each rate is
the fraction of scored units on which that obligation's measured failure
did not occur, and we define \emph{role integrity} (RI) as the
mean of the three. RI is always
reported with its components, because models lose integrity on different
obligations.
\rev{All three endpoints are failure rates, so lower is better, and their
complements in RI are the shares of units on which the measured failure was
avoided. \emph{Majority following} (OpenReview) is the share of program-chair
decisions that go with a generic majority against the paper's historical
decision; we read it against the rate at which the same model reaches that
decision from the paper alone (\S\ref{sec:or}). \emph{Error adoption}
(Court) is the share of rulings that state the planted false fact as true.
\emph{Limitation loss} (Parliament) is the share of cited witness claims that
keep neither stated limitation. RI therefore summarizes the avoidance of these
three measured failures; it does not certify that a role was fulfilled in
every other respect.}

\subsection{The three scenarios}\label{sec:scenarios}

\paragraph{OpenReview.} The area chair owes the reviewers a decision
that rests on the strength of their arguments, checked against the
paper, and not on how many of them lean each way. At ICLR
each submission receives several scored written reviews, the
authors respond in a rebuttal, an area chair writes a meta-review and a
program chair issues the decision, all publicly. Our cases are real ICLR
submissions with the paper's full text and the authors' rebuttal. In each
of our episodes, an area-chair agent reads the paper, reviews and
rebuttal and writes a meta-review that ends in a recommendation. A
program-chair agent then receives this meta-review together with the
same paper, reviews and rebuttal and issues the final decision. Every
OpenReview endpoint is computed on the program chair's decision. We
regenerate the review texts with a fixed model that reads the paper
through the same tools as the subjects, in one of two styles:
\emph{specific} reviews cite paper content, \emph{generic} reviews stay
templated. With \emph{real votes}, the submission's real 3:1 score split
is kept, the specific reviews sit with the majority (arm A) or the minority
(arm B), and the readout is whether the decision moves. With \emph{constructed votes}, the paper's
ICLR decision is the reference direction. In the control arm a specific
3:1 majority points toward it and a generic dissenter against. In the
adversarial arm a generic 3:1 majority (scores 3/3/2 or 7/7/8) points
against the reference and the lone specific dissenter toward it.
\emph{Majority following} is the share of valid adversarial-arm decisions
that go with the generic majority.

\paragraph{Court.} The judge owes counsel a ruling that checks their
assertions against the record before relying on them. In U.S.\ Supreme Court
oral argument, counsel for each side argue a case whose lower-court
history and question presented are fixed in the public record. Each of
our cases is a real case with its facts, question presented and both
sides' argument excerpts, and the senior role writes the ruling. We misstate one
procedural fact (most often the lower court's disposition) in one side's
excerpts while the record given to the judge stays correct, so the error
is checkable. Arm A places the error with the side that prevailed, arm B
with the loser, and the real outcome is withheld. Each ruling is labeled
by how it treats the planted fact: \textsc{adopted}, \textsc{corrected},
\textsc{truth stated} or \textsc{silent}. \emph{Error adoption} is how
often a ruling is \textsc{adopted}, counted over both arms.

\paragraph{Parliament.} The committee chair owes witnesses a report that
carries their evidence with the limits they placed on it. UK
parliamentary select committees run inquiries by collecting
written evidence from witnesses and publishing a report that cites it,
both public. Each of our cases is a real inquiry with its full
written evidence, and the senior role writes the report. We take one headline
empirical claim from each of two witnesses on opposite sides of the
inquiry's recorded reference and rewrite it to carry two constructed
limitations (a sample size, a scope restriction, a time window,
provisional status), checked by a second model. The readout is how the
report treats each claim: \textsc{cited with limitations}, \textsc{cited
without}, or \textsc{omitted}. \emph{Limitation loss} is how often a
cited claim has lost both limitations, counted over cited claims only.
